
%\needspace{4\baselineskip}
\color{uniblue}{\subsection{Логика отключения}}
\color{black}

\begin{enumerate}[label=\arabic{section}.\arabic{subsection}.\arabic{enumi}, labelsep=4pt, leftmargin=0pt, itemindent=57pt, itemsep=0pt, parsep=0pt] % если что можно крутить itemsep

\item
Функциональный блок логики отключения предназначен для подачи сигнала на отключение силового выключателя. 

\item
В случае присоединения линии к системе через два выключателя для каждого из них можно задействовать свою логику отключения. Далее приведено описание логики отключения для выключателя В1, логика отключения для второго выключателя аналогична.

\item
Команда на отключение выключателя \textit{«Отключить В1»} формируется:
\begin{itemize}
    \item при срабатывании внутренних защит устройства (сигнал \textit{«Ср. защит»});
    \item при срабатывании ЗНФ (сигнал \textit{«Ср. ЗНФ В1»});
    \item при действии УРОВ на отключение «своего» выключателя (сигнал \textit{«Ср. УРОВ В1 на себя»});
    \item от сигнала об аварийном давлении элегаза в ТТ, если введен программный ключ \textbf{«Откл. от Авар. давл. SF6 ТТ В1»};
    \item от команды оперативного отключения выключателя (сигнал \textit{«РКО В1»});
\end{itemize}

\item
Функциональная схема логики отключения В1 приведена на рисунке \ref{pict_logic_otkl:logic_otkl}. Функциональная схема логики отключения В2 выполняется аналогично схеме логики отключения В1.

\vspace{2mm}
\begin{figure}[H]
\centering
\includegraphics[width=1\textwidth,height=1\textheight,keepaspectratio]{\fbpath/761.9901 ЛО/img/logic_otkl.pdf}
\captionof{figure}{Функциональная схема логики отключения В1}\label{pict_logic_otkl:logic_otkl}
\end{figure}

\item
При необходимости может быть введен подхват команды на отключение от токовых реле, установленных в цепях ЭМО, в этом случае сигнал на отключение удерживается в сработанном состоянии до исчезновения тока через ЭМО (пока блок-контакты выключателя не разорвут цепи отключения). Ввод осуществляется программным ключом \textbf{«Подхват откл. В1 от РТ ЭМО»}.

\item
При наличии сигнала \textit{«Блок. отключения В1»}, формируемого функциональным блоком контроля силового выключателя, а также при выводе действия на выключатель (сигнал \textit{«Вывод действия на В1»}) отключение выключателя блокируется.

\item
Выдержка времени на возврат \textbf{«Тпрод откл. В1»} задает время продления сигнала отключения. 













\end{enumerate}